\subsection{Cantidades adoptadas para la oferta}
La revisión documental reserva un lote por cada uno de los catorce expedientes de abandono, 112 HH; la revisión escolar considera tres años de hasta 590 autorizaciones por año, 1.770 documentos: 56 lotes de hasta 32, con 8 HH por lote, 448 HH. El dimensionamiento escolar usa la proyección como capacidad de oferta, no como inventario observado. El enrolamiento reserva diez lotes de cuatro para 16 Android y 21 PC, 160 HH. La recepción de medición reserva dieciséis lotes de cuarenta puntos, 384 HH para los 640 puntos iniciales, complementando la integración del broker sin volver a presupuestarla. El CLIENTE adquiere y ejecuta las obras; Synaptix especifica, integra y recibe evidencia de calibración, precisión, asociación temporal y comportamiento del cableado.

Cada etapa reserva 107 cohortes de hasta doce asistentes: cien para usuarios finales (1.200 plazas), dos para avanzados (24), una para administradores (12), una para responsables técnicos designados (12) y tres para soporte (36). Son plazas por rol, no personas distintas: una persona que ejerce dos roles requiere las dos evaluaciones. Los 1.198 contactos potenciales del pronóstico se cubren con las 1.200 plazas finales; las categorías profesionales se añaden explícitamente, sin suponer un área TI del CLIENTE. Cada cohorte requiere 12 HH, 1.284 por etapa, más los materiales ya estimados. Las sesiones se reparten por turnos y sitios; el universo definitivo nominaliza personas y horarios dentro de esa capacidad. Se reservan 25 cohortes/300 HH en noviembre de cada uno de los tres ciclos operacionales para incorporación estacional y recuperación, adicionales a las dos jornadas anuales ya estimadas. La capacidad es ampliable por Synaptix si el universo real obligatorio lo exige y no se convierte en una exclusión de capacitación.

Se seleccionan cuatro olas por etapa, dos zonas coordinadas por ola: pantalanes/torre y zarpes; escuela/administración; varadero/combustible; portal/coordinación transversal. Durante 28 días, cada ola reserva 448 HH, 1.792 por etapa. En E1 se inicia el 1 de diciembre y toda habilitación o sesión nueva finaliza antes del día 15; desde esa fecha continúa sólo acompañamiento de los procedimientos ya activos, observación y preparación fuera de maniobras, sin capacitación o despliegue nuevo durante el congelamiento. E2 comienza en junio. La estabilización posterior reserva dos puestos diarios durante 28 días, 448 HH por etapa, adicionales a los informes técnicos de datos: en la programación seleccionada, abril de 2028 para E1 y agosto de 2028 para E2. Los cuadrantes usan catorce personas de acompañamiento y cuatro de estabilización, distintas de ingeniería y de la mesa; incluyen fines de semana y relevo.

IN2 selecciona veinte embarcaciones para el piloto, cinco lotes de revisión previa y 80 HH; su primera invernada admisible es junio de 2029 (mes 31), con inscripción y planes en mayo, tres meses de observación y botaduras en septiembre--noviembre. Los lotes se revisan al menos tres semanas antes de sus botaduras; los registros sólo dependen de los lotes aplicables ya revisados. Su evaluación y transferencia se preparan en diciembre antes del congelamiento. El seguimiento incremental 31--56 reserva 416 HH. IN1/IN4/IN5 levantan base en mayo de 2028 y evaluación en junio; IN1 imparte sus dos cohortes incrementales con dos formadores simultáneos del 29 al 31 de mayo, conservando las HH existentes. IN3 mantiene 120 pólizas y ejecución manual/asistida intercalada durante mayo--junio. Muestra insuficiente o datos inválidos impiden afirmar beneficio.

Las observaciones de recepción cuentan con ocho bloques de 80 HH, 640 HH de reserva de corrección/reensayo por etapas, asociados a sus ventanas de subsanación. No son correcciones ya realizadas ni sustituyen los diez días contractuales. Si el trabajo obligatorio excede la reserva, Synaptix incorporará capacidad a su costo y conservará los hitos y criterios de calidad. Se evita cargar el mismo cierre a la reserva y al paquete de construcción.
